%***************************************************************************************************
% config/metadata.tex - JEDYNE miejsce, które edytujesz przy tworzeniu nowego dokumentu.
% Tytuł, przedmiot, autorzy, data oraz przełączniki opcjonalnych części dokumentu.
% Nie używaj tu \today - data i skrót na stronie tytułowej są wpisywane ręcznie.
%***************************************************************************************************

% ---------------------------------------------------------------------------------------------
% Dane dokumentu (widoczne na stronie tytułowej)
% ---------------------------------------------------------------------------------------------
\newcommand{\DocType}{Dokumentacja projektowa}                       % np. Raport, Sprawozdanie, Projekt
\newcommand{\DocTitle}{Analiza poinfekcyjna artefaktów pamięci RAM oraz logów systemowych na potrzeby inżynierii detekcji}             % tytuł (może być długi, zawija się)

\newcommand{\Course}{Cyberbezpieczeństwo}        % kierunek studiów
\newcommand{\Subject}{Informatyka śledcza}             % przedmiot - pojawia się też w stopce każdej strony
\newcommand{\GroupInfo}{Grupa nr 2 / poniedziałek, godz. 11:15 - 14:00}

\newcommand{\AuthorsLabel}{Autorzy}                 % etykieta nad listą autorów (np. Autor)
\newcommand{\Authors}{%                             % linie rozdzielaj przez \\ (bez \\ na końcu)
    Aleh Shehda, nr albumu 272277\\
    Aliaksei Lapachuk, nr albumu 276711
}

\newcommand{\SupervisorLabel}{Prowadzący}           % etykieta nad danymi prowadzącego
\newcommand{\Supervisor}{%
    dr inż. Robert Czechowski
}

% Miejscowość i data - wpisywane ręcznie (świadomie bez \today).
\newcommand{\PlaceDate}{Wrocław, 2026}

% Linia ze skrótem na stronie tytułowej. Dowolny napis; pozostaw pusty (\newcommand{\DocumentHash}{}),
% aby ją ukryć. Domyślnie: skrót pustego ciągu danych (czysto dekoracyjny).
\newcommand{\DocumentHash}{4d765532542ad27ee8332d18c750238d7d4ec295ca117fb7aeca216a0339bdf0}

% Autor w metadanych pliku PDF (bez znaków \\ i adresów e-mail)
\newcommand{\PdfAuthor}{Aleh Shehda, Aliaksei Lapachuk}

% ---------------------------------------------------------------------------------------------
% Przełączniki opcjonalnych części dokumentu: \...true = pokaż, \...false = ukryj
% ---------------------------------------------------------------------------------------------
\newif\ifshowabstract         \showabstractfalse          % Streszczenie przed spisem treści
\newif\ifshowabbreviations    \showabbreviationstrue      % Spis skrótów po spisie treści
\newif\ifshowbibliography     \showbibliographytrue       % Bibliografia (BibTeX, bib/references.bib)
\newif\ifshowlistoffigures    \showlistoffigurestrue      % Spis rysunków
\newif\ifshowlistoftables     \showlistoftablestrue       % Spis tabel
\newif\ifshowlistoflistings   \showlistoflistingstrue     % Spis listingów
\newif\ifshowlistofalgorithms \showlistofalgorithmsfalse  % Spis algorytmów
\newif\ifshowappendices       \showappendicesfalse        % Załączniki

% Wersja do druku: wyłącza kolorowe odnośniki w PDF (\printversiontrue)
\newif\ifprintversion         \printversionfalse

% Domyślne umiejscowienie rysunków wstawianych makrem \Figure
% (H = dokładnie w miejscu wywołania; htbp = pozwól LaTeX-owi przesuwać)
\newcommand{\DefaultFigurePlacement}{H}

